\section{Gradient bounds for concave equations}
\label{sec:concave-gradient}\index{Gradient estimate!concave equations}

Keep the fixed data and normalization of \cref{sec:concave-strict-gap-stability}, and assume in addition that $\|\partial h\|_{L^\infty(\omega)}\leq H_1$. The stability estimate supplies nearby smooth admissible functions. A weighted gradient comparison with one such function gives a bound whose comparison radius is fixed before its derivatives are known. Compactness then reduces the estimate to finitely many smooth comparison functions.

For the derivative calculation, write $U=u/(2\pi)$ and $\underline U=\underline u/(2\pi)$. Thus $\ddc u=\mathrm i\partial\bar\partial U$. For any smooth real $V$, put $\mathcal A_V=\omega^{-1}(\chi+\mathrm i\partial\bar\partial V)$ and call $V$ admissible if $\lambda(\mathcal A_V)\in\Gamma$. Use the Chern connection of $\omega$ and define $B=DF(\mathcal A_U)$ by the trace pairing. At a point where $\mathcal A_U$ is diagonal in a unitary frame, set
\[
 B_i=f_i(\lambda(\mathcal A_U)),\qquad
 \mathcal Lq=\sum_iB_iq_{i\bar i},\qquad T=\sum_iB_i,
 \qquad |\partial q|_B^2=\sum_iB_i|q_i|^2.
\]
Here $U_{ki}=\nabla_iU_k$ and $U_{i\bar j}=\partial_i\partial_{\bar j}U$; the pure derivatives need not be symmetric. The trace expressions below are invariant under a change of unitary frame.

\Needspace{11\baselineskip}
\begin{lemma}\label{lem:concave-linearized-alternative}
There are fixed $\tau,\kappa>0$ such that
\begin{gather*}
 T\geq\tau,\qquad 0\leq\operatorname{tr}(B\mathcal A_U)\leq h,\\
 \mathcal L(U-V)\leq h\quad(V\text{ admissible}),\qquad
 \mathcal L(\underline U-U)\geq-h.
\end{gather*}
At every point, at least one of the alternatives
\[
 \mathcal L(\underline U-U)\geq\kappa T,
 \qquad B\geq\kappa T\operatorname{Id}
\]
holds. The constants are uniform for all positive right-hand sides in the family.
\end{lemma}
\begin{proof}
If $P$ is admissible, the ray $\mathcal A_U+tP$ is admissible. Its positive concave operator value cannot have a negative tangent slope, so $\operatorname{tr}(BP)\geq0$. The tangent inequality at the origin, where $F(0)=0$, gives $\operatorname{tr}(B\mathcal A_U)\leq h$. Choose $a_*>0$ with $F(a_*\operatorname{Id})>H_0+1$. Concavity gives
\[
 H_0+1<F(a_*\operatorname{Id})
 \leq h+a_*T-\operatorname{tr}(B\mathcal A_U)\leq H_0+a_*T.
\]
Thus $T\geq a_*^{-1}$. The remaining two bounds follow by subtracting the nonnegative traces of $\mathcal A_V$ and $\mathcal A_{\underline U}$.

We prove the alternative directly from the strict gap. Fix $r,H$ from \cref{lem:strict-gap-matrix-bound}, and put $\mu=\lambda(\mathcal A_{\underline U})$. Compactness and the gap give one $S>r$ such that
\[
 f(\mu-2r\mathbf1+Se_i)>H+2\delta\qquad(1\leq i\leq n)
\]
everywhere. For $\lambda\in\Gamma$ with $f(\lambda)=\sigma=h$, consider the compact set
\[
 K=\{v\geq\mu-r\mathbf1:f(v)\leq\sigma\}.
\]
It is contained in $\Gamma$ because $\mu-r\mathbf1\in\Gamma$; the matrix bound makes it bounded. When nonempty, it contains $\mu-r\mathbf1$ and is star-shaped about that point: every segment towards it decreases each coordinate, hence decreases $f$. Thus it is connected. Let $N=\nabla f(\lambda)$, whose entries are positive, and let $\Pi$ be the tangent hyperplane $N\cdot(v-\lambda)=0$.

If $f(\mu-r\mathbf1)>\sigma$, concavity gives $N\cdot(\mu-\lambda)>r\sum_iN_i$. Otherwise $\mu-r\mathbf1\in K$. The positive diagonal ray from that point meets $\{f=\sigma\}$, since $f$ diverges along this ray. At the meeting point the tangent inequality gives $N\cdot(v-\lambda)\geq0$. If $\Pi$ misses $K$, connectedness makes this sign strictly positive throughout $K$, again giving $N\cdot(\mu-\lambda)>r\sum_iN_i$.

If $v\in\Pi\cap K$, put $z_i=v-r\mathbf1+Se_i$. Then $z_i\geq\mu-2r\mathbf1+Se_i$, so $f(z_i)>\sigma$. The tangent inequality at $\lambda$ gives
\[
 0<N\cdot(z_i-v)=SN_i-r\sum_jN_j.
\]
Thus $N_i>(r/S)\sum_jN_j$ for every $i$. We have proved the required scalar alternative at every positive level, with constants independent of $\sigma$.

Order $\lambda$ and $\mu$ decreasingly. Concavity and symmetry order $f_i(\lambda)$ increasingly. In the eigenframe of $\mathcal A_U$, the diagonal of $\mathcal A_{\underline U}$ is a doubly stochastic average of permutations of $\mu$. The rearrangement inequality therefore gives
\[
 \operatorname{tr}\bigl(B(\mathcal A_{\underline U}-\mathcal A_U)\bigr)
 \geq\sum_if_i(\lambda)(\mu_i-\lambda_i).
\]
Taking $\kappa$ smaller than $\min\{r,r/S\}$ proves the matrix alternative.
\end{proof}

\Needspace{10\baselineskip}
\begin{lemma}\label{lem:concave-weighted-gradient}
Put $Q=1+|\partial U|_\omega^2$. For any smooth real $\xi$, one has
\begin{align*}
 \mathrm e^{-\xi}\mathcal L(\mathrm e^\xi Q)
 &\geq\sum_{i,k}B_i|U_{ki}+U_k\xi_i|^2
       +\frac14\sum_iB_i\lambda_i^2+Q\mathcal L\xi+|\partial\xi|_B^2\\
 &\quad+2\Real\sum_iB_i\lambda_iU_i\xi_{\bar i}
       -2\Real\sum_{i,k}B_i\chi_{i\bar k}U_k\xi_{\bar i}
       -C_gTQ-2H_1|\partial U|_\omega,
\end{align*}
where $C_g$ depends only on $\omega$ and $\chi$.
\end{lemma}
\begin{proof}
Write $\mathscr T^p_{ik}$ and $R^p_{i\bar jk}$ for Chern torsion and curvature. The commutation and differentiated equation are
\begin{gather*}
 \nabla_{\bar j}\nabla_iU_k
 =\nabla_kU_{i\bar j}+\sum_p\bigl(R^p_{i\bar jk}U_p-\mathscr T^p_{ik}U_{p\bar j}\bigr),\\
 \sum_iB_i\nabla_kU_{i\bar i}=h_k-\sum_iB_i\nabla_k\chi_{i\bar i}.
\end{gather*}
The product rule for $|\partial U|^2$ gives both squares $\sum B_i|U_{ki}|^2$ and $\sum B_i|U_{k\bar i}|^2$, the term $2\Real\sum_kh_kU_{\bar k}$, a curvature error bounded by $CT|\partial U|^2$, and the errors from $\nabla\chi$ and torsion. The latter are bounded by
\[
 CT|\partial U|+CT|\partial U|^2
 +C|\partial U|T^{1/2}\Bigl(\sum_{i,k}B_i|U_{k\bar i}|^2\Bigr)^{1/2}.
\]
Absorb half the mixed square. Since $U_{i\bar j}=\lambda_i\delta_{ij}-\chi_{i\bar j}$, the retained half bounds $\frac14\sum_iB_i\lambda_i^2-CT$. Hence
\[
 \mathcal L|\partial U|^2\geq\sum_{i,k}B_i|U_{ki}|^2
 +\frac14\sum_iB_i\lambda_i^2-C_gTQ-2H_1|\partial U|.
\]
Finally,
\[
 \partial_i|\partial U|^2
 =\sum_kU_{ki}U_{\bar k}+\lambda_iU_i-\sum_k\chi_{i\bar k}U_k.
\]
In the product rule for $\mathrm e^\xi Q$, combine its first sum with $|\partial U|^2|\partial\xi|_B^2$ and the pure-derivative square. This gives $\sum B_i|U_{ki}+U_k\xi_i|^2$, while the remaining weight term is $|\partial\xi|_B^2$.
\end{proof}

\Needspace{10\baselineskip}
\begin{proposition}\label{prop:concave-nearby-gradient}
There are fixed $\eta,C>0$ such that every smooth admissible $V$ with $\|U-V\|_\infty\leq\eta$ satisfies
\[
 \|\partial U\|_\infty\leq C\bigl(1+\|\partial V\|_\infty^2\bigr).
\]
The comparison radius $\eta$ is independent of the derivatives of $V$.
\end{proposition}
\begin{proof}
Let $a\geq1$ and $\beta>1$ be chosen below. Assuming $\|U-V\|_\infty\leq\beta^{-1}$, put
\[
 w=U-V-\beta^{-1},\qquad
 \xi=a(\underline U-U)+\beta w^2,\qquad
 K_V=1+\|\partial\underline U\|_\infty+\|\partial V\|_\infty.
\]
Then $-2\beta^{-1}\leq w\leq0$ and
\[
 \xi_i=a(\underline U_i-V_i)+(-a+2\beta w)(U_i-V_i),
 \qquad |-a+2\beta w|\leq a+4.
\]
Writing $D=|\partial(U-V)|_B^2$, the Cauchy--Schwarz inequality gives
\begin{align*}
 \left|2\Real\sum_iB_i\lambda_iU_i\xi_{\bar i}\right|
 &\leq\frac18\sum_iB_i\lambda_i^2+C_aK_V^4T+C_a(Q+K_V^2)D,\\
 \left|2\Real\sum_{i,k}B_i\chi_{i\bar k}U_k\xi_{\bar i}\right|
 &\leq\tfrac12QD+C_aT+C_aK_VQ^{1/2}T.
\end{align*}
For the first line, split $\xi_i$ as above and use $\sum B_i|U_i|^2\leq2D+2K_V^2T$. Apply Young's inequality to the resulting three products, each absorbing at most $\frac1{24}\sum B_i\lambda_i^2$. For the second line use $|\partial\xi|_B\leq aK_VT^{1/2}+(a+4)D^{1/2}$ and the fixed bound for $\chi$. Thus $C_a$ is independent of $\beta$ and of the derivatives of $V$.

By \cref{lem:concave-linearized-alternative} and $w\leq0$,
\[
 \mathcal L\xi\geq a\mathcal L(\underline U-U)-4H_0+2\beta D.
\]
Choose $a$ large enough that $(a\kappa-C_g-2)\tau>4H_0+2H_1+1$, and then $\beta\geq2C_a+2$. Combining the preceding bounds with \cref{lem:concave-weighted-gradient}, whenever $Q\geq C'(1+K_V^4)$, gives
\begin{align*}
 \mathrm e^{-\xi}\mathcal L(\mathrm e^\xi Q)
 &\geq\frac18\sum_iB_i\lambda_i^2\\
 &\quad+Q\bigl[a\mathcal L(\underline U-U)-(C_g+2)T
                 +\beta D-(4H_0+2H_1)\bigr].
\end{align*}
Indeed, the coefficient of $D$ before absorption is $(2\beta-C_a-1/2)Q-C_aK_V^2$, which is at least $\beta Q$ after increasing the fixed $C'$. The other errors satisfy $C_aK_V^4\leq Q$ and $C_a+C_aK_VQ^{1/2}\leq Q$, giving the additional $2T$.

At a maximum of $\mathrm e^\xi Q$, suppose this large-gradient inequality holds. In the first linearized alternative the bracket is positive by the choice of $a$. In the second, $D\geq\kappa T|\partial(U-V)|^2$. For $Q\geq4(1+K_V^2)$, the triangle inequality gives $|\partial(U-V)|^2\geq c_*Q$, where $c_*=(\sqrt3-1)^2/4$. Since $\mathcal L(\underline U-U)\geq-H_0$, the bracket is at least
\[
 (\beta\kappa c_*Q-C_g-2)T-(a+4)H_0-2H_1.
\]
It is positive once the fixed $C'$ is large enough, because $T\geq\tau$. Both cases contradict the second derivative test. Thus $Q$ at the maximum is bounded by $C(1+K_V^4)$. The weight has uniformly bounded oscillation:
\[
 \operatorname{osc}_X\xi\leq a(\operatorname{osc}_X\underline U+\operatorname{osc}_XU)+4/\beta.
\]
The potential bound in \cref{prop:strict-gap-potential-bound} therefore transfers the maximum-point estimate to every point. Taking square roots and setting $\eta=\beta^{-1}$ proves the assertion.
\end{proof}

\Needspace{9\baselineskip}
\begin{theorem}\label{thm:concave-uniform-gradient}
Every solution with the fixed data and $\|\partial h\|_\infty\leq H_1$ satisfies $\|\mathrm du\|_{L^\infty(\omega)}\leq C$. The constant is uniform as the positive right-hand sides tend to zero.
\end{theorem}
\begin{proof}
The entire family of right-hand sides with the fixed upper and first-derivative bounds is equicontinuous. By \cref{cor:concave-uniform-compactness}, the corresponding normalized solutions $U$ form a precompact family in $C^0(X)$. Choose a finite $\eta$-net from this family, with smooth admissible members $V_1,\ldots,V_N$. Each solution is within $\eta$ of one $V_j$, so \cref{prop:concave-nearby-gradient} bounds its gradient by $C(1+\max_j\|\partial V_j\|_\infty^2)$. The full family and the chosen finite net depend only on the fixed data and $H_1$. Finally $u=2\pi U$, so the same conclusion holds for $\mathrm du$.
\end{proof}
